\chapter{Results}
\label{chap:results}

\section{Overview}

The final results show that checkpoint loading performance is strongly workload-dependent. The clearest practical framing is an adaptive heuristic: use a simple synchronous path for smaller eager loads, and switch to io\_uring once checkpoint size and shard structure become large enough to justify a deeper submission path.

At the Rust layer, \texttt{SafeTensors} and \texttt{ServerlessLLM} show different backend preferences. At the Python layer, binding overhead and tensor materialisation partially obscure raw Rust backend gains. At the vLLM layer, the effect of backend choice is strongest for model initialisation and time-to-first-token, while steady-state decode is comparatively insensitive.

These results support the main thesis of this project: the practical goal is to build a workload-aware heuristic for LLM checkpoint loading, not to prove that a single backend should replace all others.

\section{Rust-Layer Backend Results}

The Rust harness provides the cleanest backend-attribution layer because it measures checkpoint loading directly, without Python tensor conversion or serving-engine setup costs. The full cold-cache matrix is stored in \texttt{results/h100/profile/full\_cold\_matrix.tsv}.

\subsection{SafeTensors}

Table~\ref{tab:rust-safetensors} summarises the explicit backend comparison for \texttt{SafeTensors}. Two regimes are clear.

\begin{table}[h]
\caption{Rust cold-cache \texttt{SafeTensors} results (ms). Best explicit backend highlighted in prose.}
\label{tab:rust-safetensors}
\centering
\begin{tabular}{lrrrr}
\toprule
Model & Sync & Async & Default & io\_uring \\
\midrule
Qwen3-0.6B & 555.28 & 1416.19 & 333.59 & 1159.24 \\
SmolLM3-3B & 1331.79 & 5802.05 & 3722.35 & 3812.35 \\
Qwen3-4B & 1820.85 & 5039.21 & 3083.18 & 2967.42 \\
Qwen3-8B & 3880.40 & 7074.40 & 3844.35 & 3172.37 \\
Qwen3-14B & 7030.25 & 13489.00 & 6138.70 & 5977.07 \\
Qwen3-32B & 15454.45 & 28763.46 & 12119.50 & 12089.05 \\
\bottomrule
\end{tabular}
\end{table}

For smaller eager whole-file loads, synchronous loading remains highly effective. On SmolLM3-3B and Qwen3-4B, \texttt{sync} is the strongest explicit backend. However, the ranking flips on larger multi-shard checkpoints. On Qwen3-8B, Qwen3-14B, and Qwen3-32B, the multi-worker \texttt{io\_uring} backend is the fastest explicit backend. The resulting rule is the clearest high-level heuristic in the project: \texttt{sync} for smaller \texttt{SafeTensors} checkpoints, \texttt{io\_uring} for larger ones.

The \texttt{async} backend is consistently the weakest choice for \texttt{SafeTensors}. This matches the workload structure: whole-file shard loading benefits from strong bulk parallelism, not from scheduler-heavy fine-grained async orchestration.

Repeated five-run anchor measurements reinforce this regime split. For the anchor set, the mean best backend is \texttt{sync} for Qwen3-0.6B and \texttt{io\_uring} for Qwen3-8B and Qwen3-32B.

\subsection{ServerlessLLM}

The \texttt{ServerlessLLM} layout produces a different ordering, shown in Table~\ref{tab:rust-serverlessllm}.

\begin{table}[h]
\caption{Rust cold-cache \texttt{ServerlessLLM} results (ms).}
\label{tab:rust-serverlessllm}
\centering
\begin{tabular}{lrrrr}
\toprule
Model & Sync & Async & Default & io\_uring \\
\midrule
Qwen3-0.6B & 1645.06 & 1471.34 & 1432.25 & 1387.57 \\
SmolLM3-3B & 3941.88 & 2391.24 & 3182.28 & 3091.49 \\
Qwen3-4B & 5137.72 & 3504.25 & 3551.82 & 3727.14 \\
Qwen3-8B & 11048.40 & 9463.37 & 7589.47 & 9043.51 \\
Qwen3-14B & 18661.25 & 16139.08 & 13422.94 & 13257.23 \\
Qwen3-32B & 39822.09 & 29180.89 & 27659.78 & 29779.86 \\
\bottomrule
\end{tabular}
\end{table}

Here, \texttt{sync} is never the best backend. For smaller and mid-sized range-heavy workloads, \texttt{async} is often strong because grouped and coalesced range reads fit its execution model well. For larger cases, \texttt{io\_uring} becomes competitive or best. For example, Qwen3-14B is fastest with \texttt{io\_uring}, while SmolLM3-3B and Qwen3-4B favour \texttt{async}.

The key result is not that one backend dominates, but that the partitioned layout changes the ranking. This is exactly the kind of workload sensitivity that motivates an adaptive \texttt{default} policy built around workload-aware heuristics rather than fixed backend dogma.

\subsection{Selector Behaviour}

The final selector is not perfect on every single Rust matrix row, but it is materially better than the earlier cost-model-based selector. On the Rust matrix, \texttt{default} is exactly best for Qwen3-8B and Qwen3-32B in \texttt{ServerlessLLM}, and very close to best for the large \texttt{SafeTensors} cases. The remaining misses are concentrated around the medium-size crossover region, especially SmolLM3-3B and Qwen3-4B.

The important point is that the selector now follows the correct \emph{regimes}: smaller \texttt{SafeTensors} loads lean toward \texttt{sync}; larger \texttt{SafeTensors} loads lean toward \texttt{io\_uring}; \texttt{ServerlessLLM} switches between \texttt{async} and \texttt{io\_uring} depending on workload size and structure. In other words, the final selector behaves like a practical checkpoint-loading heuristic rather than a static backend preference.

\section{Python Binding Results}

Python benchmarks were collected on Qwen/Qwen3-8B to quantify user-facing overhead and to evaluate whether binding optimisation and selector redesign improved practical performance.

\subsection{Binding Overhead Before and After Optimisation}

The most important Python result is that binding engineering mattered. Replacing per-call runtime setup with a shared Tokio runtime and reducing repeated Torch conversion work materially improved the public \texttt{default} API.

For \texttt{SafeTensors} on Qwen3-8B:

\begin{itemize}
\item pre-optimisation \texttt{default}: 15.8831~s
\item post-binding-optimisation \texttt{default}: 14.5451~s
\item post-selector-redesign \texttt{default}: 11.5498~s
\end{itemize}

This is a 20.6\% reduction between the post-optimisation and post-selector runs, and a much larger improvement relative to the earliest Python baseline.

For \texttt{ServerlessLLM} on Qwen3-8B:

\begin{itemize}
\item early \texttt{default}: 22.4547~s
\item post-binding-optimisation \texttt{default}: 20.5008~s
\item post-selector-redesign \texttt{default}: 18.4518~s
\end{itemize}

This is a 10.0\% reduction between the post-optimisation and post-selector runs.

\subsection{Python Layer Interpretation}

The Python numbers are much larger than the Rust-only numbers because eager loading at the binding layer includes substantial tensor materialisation and Python object construction. This is especially visible in the \texttt{SafeTensors} benchmark, where native \texttt{safetensors} loading remains much faster than the fully materialised \texttt{tensor\_store} eager paths.

The key conclusion is that Python is not a transparent pass-through layer. Backend improvements matter, but they are filtered through the cost of converting checkpoint contents into Torch tensors and Python dictionaries. That is why the binding optimisations were necessary before the Python layer could reflect the improved Rust backends more faithfully.

\section{vLLM Results}

The final vLLM matrix completed successfully across all six models, all benchmark loaders, warm and cold cache conditions, and all three benchmark kinds. The completed summary is stored in \texttt{results/h100/vllm/full\_matrix.tsv}. The final matrix contains 672 successful rows and no failures.

\subsection{Load-Only Results}

For cold \texttt{load\_only}, the best loader varies with model size. Table~\ref{tab:vllm-load-only} shows the strongest loader in the final cold-load-only configuration.

\begin{table}[h]
\caption{Best cold \texttt{load\_only} vLLM loader by model.}
\label{tab:vllm-load-only}
\centering
\begin{tabular}{lrl}
\toprule
Model & Best init time (ms) & Best loader \\
\midrule
Qwen3-0.6B & 20417.42 & ts\_safetensors\_sync \\
SmolLM3-3B & 27333.96 & ts\_safetensors\_default \\
Qwen3-4B & 28475.27 & ts\_safetensors\_default \\
Qwen3-8B & 35555.79 & ts\_safetensors\_default \\
Qwen3-14B & 47751.38 & ts\_safetensors\_io\_uring \\
Qwen3-32B & 88164.00 & ts\_safetensors\_io\_uring \\
\bottomrule
\end{tabular}
\end{table}

Two observations stand out. First, the best vLLM loader is almost always from the \texttt{SafeTensors} family rather than the \texttt{ServerlessLLM} family. Second, the explicit \texttt{io\_uring} path becomes valuable again at the larger end of the ladder. It is the best \texttt{load\_only} choice for Qwen3-14B and Qwen3-32B.

\subsection{TTFT and Steady-State Decode}

The backend effect is strongest for initialisation and TTFT, and much weaker for steady-state decode. This is exactly what one would expect: the loading backend can change how quickly weights become resident, but once the model is fully loaded, decode throughput is dominated by the GPU compute path.

For cold TTFT, \texttt{ts\_safetensors\_default} is the strongest loader for Qwen3-4B, Qwen3-8B, and Qwen3-14B, while \texttt{ts\_safetensors\_sync} remains best for Qwen3-0.6B and SmolLM3-3B. For steady-state decode, differences are much smaller and the winning loader varies more irregularly. This confirms that backend choice is primarily a cold-start and initialisation problem rather than a decode-throughput problem.

\subsection{Large-Model Feasibility}

The 14B and 32B models initially failed in vLLM. These were not backend failures. The raw logs showed that the benchmark harness was too conservative in its vLLM memory budget, leaving insufficient KV-cache allocation headroom. After increasing the vLLM memory budget and reducing \texttt{max\_model\_len}, all previously failing rows succeeded.

This matters for interpretation. The initial failures were a serving-harness configuration issue, not evidence that the checkpoint backends could not support the large models.

\section{Negative and Turnaround Results}

The most important negative result in the project was the early single-ring \texttt{io\_uring} prototype. On the H100 system, that implementation was much slower than \texttt{sync} and showed poor CPU utilisation. Profiling pointed to a design problem rather than a fundamental backend limitation: one ring and one issuer thread left too much of the machine idle.

The project changed direction after that result. Once the \texttt{io\_uring} backend was redesigned around multiple worker-owned rings, its ranking changed significantly. It became competitive or best on larger \texttt{SafeTensors} and \texttt{ServerlessLLM} workloads. The final result is therefore not ``io\_uring always wins'', nor ``io\_uring is a dead end''. The correct conclusion is that backend family and backend architecture both matter.

\section{Summary}

Three final findings emerge from the completed result set.

\begin{enumerate}
\item At the Rust layer, backend ranking depends strongly on workload structure. \texttt{SafeTensors} favours \texttt{sync} on smaller cases and \texttt{io\_uring} on larger cases. \texttt{ServerlessLLM} favours \texttt{async} or \texttt{io\_uring}, with \texttt{sync} rarely best.
\item At the Python layer, binding overhead is large enough to matter, and binding-specific optimisation produced meaningful practical gains.
\item At the vLLM layer, backend choice affects initialisation and TTFT far more than steady-state decode, and the best loader depends on model size.
\end{enumerate}

Taken together, the results support an adaptive, format-aware checkpoint-loading heuristic. The dominant practical rule is sync for smaller models and io\_uring for larger ones, with additional format-specific adjustments for range-heavy layouts.
